\documentclass[uplatex,a4paper,12pt]{jsarticle}
\usepackage{amsmath,amssymb}
\usepackage{enumitem}

\begin{document}

\title{集合・微分 確認テスト}
\author{}
\date{}
\maketitle

\section*{注意}
\begin{itemize}
\item 試験時間は5分とする．
\item 全10問，各10点，合計100点とする．
\item 60点以上を合格とする．
\end{itemize}

\section*{問題}

\begin{enumerate}

\item 集合の説明として正しいものを1つ選べ．
\begin{enumerate}[label=(\alph*)]
\item ある条件を満たすものの集まり
\item 1つの数だけを表すもの
\item 微分するための記号
\end{enumerate}

\item $a$ が集合 $A$ の要素であることを表す記号として正しいものを1つ選べ．
\begin{enumerate}[label=(\alph*)]
\item $a \in A$
\item $a = A$
\item $a \to A$
\end{enumerate}

\item 写像
\[
f:A\to B
\]
の説明として正しいものを1つ選べ．
\begin{enumerate}[label=(\alph*)]
\item $A$ の元を $B$ の元に対応させる規則
\item $A$ と $B$ を足し算する規則
\item $A$ の元をすべて削除する規則
\end{enumerate}

\item 集合
\[
A=\{x\in \mathbb{R}\mid 1<x<5\}
\]
について，上限 $\sup A$ と下限 $\inf A$ を求めよ．

\item 次の関数を $x$ について偏微分せよ．
\[
f(x,y)=3x^2+2xy+y^2
\]

\item 次の関数を $y$ について偏微分せよ．
\[
f(x,y)=3x^2+2xy+y^2
\]

\item 次の関数の勾配 $\nabla f$ を求めよ．
\[
f(x,y)=x^2+4y^2
\]

\item 次の変数変換について，ヤコビ行列 $J$ を求めよ．
\[
u=x+y,\quad v=x-y
\]

\item 次の変数変換について，ヤコビアン $|J|$ を求めよ．
\[
u=x+y,\quad v=x-y
\]

\item
\[
x=uv,\quad y=u+v,\quad f(x,y)=x^2+y^2
\]
とする．連鎖律を用いるとき，
\[
\frac{\partial f}{\partial u}
\]
として正しいものを1つ選べ．
\begin{enumerate}[label=(\alph*)]
\item $2uv^2+2u+2v$
\item $2u^2v+2u+2v$
\item $uv+u+v$
\end{enumerate}

\end{enumerate}

\end{document}